\begin{table}[htbp]
\centering
\footnotesize
\begin{threeparttable}
\caption{Resource ontology used throughout the manuscript. The supplementary atlas maps these reader-facing labels to normalized repository evidence identifiers.}
\label{tab:resource-ontology}
\begin{tabularx}{\linewidth}{@{}p{0.08\linewidth}p{0.25\linewidth}X@{}}
\toprule
\textbf{Tag} & \textbf{Added resource} & \textbf{Interpretation} \\
\midrule
R1 & More coordinates & Larger representational capacity or a wider margin budget. \\
R2 & More coordinate precision & More bits or analog state per coordinate, preserving soft evidence that binary views discard. \\
R3 & Sparse or block structure & Reduced interference or more structured evidence placement across coordinates. \\
R4 & Native code structure & Explicit algebraic or combinatorial constraints rather than random assignment alone. \\
R5 & Exact side information & Typed exact metadata, handles, or manifests that preserve authority outside the lossy semantic view. \\
R6 & External context or prior & Candidate-space narrowing before the main decoder spends its full search budget. \\
R7 & Decoder compute & More iterations, stronger search, or a more expensive inference procedure. \\
R8 & Restarts or wider search & Additional exploration budget that tries more trajectories or candidates. \\
R9 & Temporal state & More time steps or internal state evolution during recovery. \\
R10 & Physical parallelism & Device-level concurrency or local state that changes the practical price of recovery. \\
R11 & Preprocessing or materialization & Upfront indexing, cached views, or extra write-time work. \\
R12 & Dual representation & Persistent maintenance of multiple native views for the same underlying state. \\
R13 & Reduced coverage through abstention & Refusal to answer under insufficient evidence instead of forcing acceptance. \\
R14 & Exact fallback & Retaining or invoking a non-lossy recovery path when the approximate route is insufficient. \\
\bottomrule
\end{tabularx}
\end{threeparttable}
\end{table}
